% Lecture 2 notes: \input by notes.tex, \subimport'ed by the semester book.
% Exam questions (экзамен): 3 (DLP в группе порядка q^e: сведение к подгруппе порядка q),
% 4 (китайская теорема об остатках; алгоритм Полига–Хеллмана, оценка сложности).
\begin{lecture}{2}{Китайская теорема об остатках и алгоритм Полига–Хеллмана}

\section{Почему важен порядок группы}

В лекции~1 мы познакомились с задачей дискретного логарифмирования (DLP): в циклической группе
$G = \langle g \rangle$ порядка $N$ по элементу $h \in G$ требуется найти показатель $x$, для
которого $g^x = h$. Протокол Диффи"--~Хеллмана стоек лишь постольку, поскольку эта задача сложна;
универсальный алгоритм Шенкса решает её за $\bigO(\sqrt N)$ групповых операций.

Оказывается, сложность DLP определяется не самим числом $N$, а его \emph{разложением на простые
множители}. Если $N = N_1 N_2$ с взаимно простыми $N_1$, $N_2$, то задача распадается на две
задачи в группах порядков $N_1$ и $N_2$, ответы которых «склеиваются» китайской теоремой об
остатках. Если же $N = q^e$ "--- степень простого числа, то $x$ можно находить по одной цифре
в $q$-ичной системе, решая каждый раз задачу в группе порядка всего лишь $q$. Вместе эти два
наблюдения составляют \emph{алгоритм Полига"--~Хеллмана}~\cite{pohlig-hellman}: он решает DLP
за время, определяемое \emph{наибольшим простым делителем} порядка группы.

Пример, который стоит держать в голове: $p = 2^{16} + 1 = 65537$ "--- простое число Ферма, и
группа $\Znx{p}$ имеет порядок $2^{16}$. Метод Шенкса потребовал бы около $2^8$ операций, а
алгоритм этой лекции находит логарифм по одному биту, $16$ раз сравнивая элемент с $\pm 1$.
Для криптографии вывод однозначен: порядок рабочей группы обязан иметь большой простой делитель.

\section{Китайская теорема об остатках}

Напомним (лекция~3 первого семестра), что $a$ обратим по модулю $m$ тогда и только тогда,
когда $\gcd(a, m) = 1$, а обратный элемент находится расширенным алгоритмом Евклида: из
соотношения Безу $ua + vm = 1$ получаем $a^{-1} \equiv u \pmod m$.

\begin{lemma}\label{l2:coprime-div}
Если $a \mid n$, $b \mid n$ и $\gcd(a, b) = 1$, то $ab \mid n$. Как следствие, если
$m_1, \dots, m_k$ попарно взаимно просты и каждое $m_i$ делит $n$, то $m_1 \cdots m_k \mid n$.
\end{lemma}

\begin{proof}
Пусть $n = a n'$ и $ua + vb = 1$. Тогда $n' = u a n' + v b n' = u n + v b n'$, и оба слагаемых
делятся на $b$ (так как $b \mid n$). Значит, $b \mid n'$, т.\,е. $ab \mid n$. Второе утверждение
"--- индукция по $k$: число $m_k$ взаимно просто с произведением $m_1 \cdots m_{k-1}$ (у них нет
общих простых делителей).
\end{proof}

\begin{theorem}[китайская теорема об остатках, КТО]\label{l2:crt}
Пусть $m_1, \dots, m_k$ "--- попарно взаимно простые натуральные числа, $M = m_1 m_2 \cdots m_k$.
Тогда для любых целых $r_1, \dots, r_k$ система
\[
  x \equiv r_1 \Mod{m_1}, \quad x \equiv r_2 \Mod{m_2}, \quad \dots, \quad x \equiv r_k \Mod{m_k}
\]
имеет решение, и оно единственно по модулю $M$. Явно:
\[
  x \equiv \sum_{i=1}^{k} r_i M_i y_i \pmod M, \qquad M_i = \frac{M}{m_i}, \quad
  y_i \equiv M_i^{-1} \pmod{m_i}.
\]
\end{theorem}

\begin{proof}
\emph{Существование.} Число $M_i$ "--- произведение модулей, отличных от $m_i$, поэтому
$\gcd(M_i, m_i) = 1$ и обратный $y_i$ существует. При $j \ne i$ имеем $m_j \mid M_i$, следовательно,
\[
  M_i y_i \equiv \begin{cases} 1 \pmod{m_i},\\ 0 \pmod{m_j},\ j \ne i.\end{cases}
\]
Поэтому для $x = \sum_i r_i M_i y_i$ по модулю $m_j$ остаётся единственное слагаемое:
$x \equiv r_j M_j y_j \equiv r_j \pmod{m_j}$.

\emph{Единственность.} Если $x$ и $x'$ "--- два решения, то $m_i \mid x - x'$ при всех $i$,
и по лемме~\ref{l2:coprime-div} $M \mid x - x'$. Наоборот, любое $x' \equiv x \pmod M$ также
является решением.
\end{proof}

Числа $e_i = M_i y_i$ (\emph{идемпотенты}: $e_i^2 \equiv e_i$, $e_i e_j \equiv 0 \pmod M$)
зависят только от модулей; если система решается много раз с одними и теми же модулями, их
вычисляют один раз, и решение сводится к $k$ умножениям и сложениям.

\begin{example}\label{l2:suntzu}
Задача из трактата «Сунь"--~цзы суань цзин» (III"--~V вв.): найти число, которое при делении
на $3$ даёт остаток $2$, на $5$ "--- остаток $3$, на $7$ "--- остаток $2$. Здесь $M = 105$,
\[
\begin{gathered}
  M_1 = 35 \equiv 2 \pmod 3,\ y_1 = 2; \qquad M_2 = 21 \equiv 1 \pmod 5,\ y_2 = 1;\\
  M_3 = 15 \equiv 1 \pmod 7,\ y_3 = 1,
\end{gathered}
\]
\[
  x \equiv 2 \cdot 35 \cdot 2 + 3 \cdot 21 \cdot 1 + 2 \cdot 15 \cdot 1 = 233 \equiv 23 \pmod{105}.
\]
\end{example}

\taskexample{crt-coprime}

\subsection{КТО как изоморфизм колец}

\begin{theorem}\label{l2:crt-iso}
Если $m_1, \dots, m_k$ попарно взаимно просты и $M = m_1 \cdots m_k$, то отображение
\[
  \varphi\colon \Zn{M} \to \Zn{m_1} \times \cdots \times \Zn{m_k}, \qquad
  x \bmod M \mapsto (x \bmod m_1, \dots, x \bmod m_k),
\]
"--- изоморфизм колец. В частности,
$\Znx{M} \cong \Znx{m_1} \times \cdots \times \Znx{m_k}$ и функция Эйлера мультипликативна:
$\varphi(M) = \varphi(m_1) \cdots \varphi(m_k)$.
\end{theorem}

\begin{proof}
Отображение определено корректно: если $x \equiv x' \pmod M$, то $x \equiv x' \pmod{m_i}$, так как
$m_i \mid M$. Оно сохраняет сложение и умножение, поскольку их сохраняет каждая проекция
$\ZZ \to \Zn{m_i}$. Сюръективность "--- это существование решения в теореме~\ref{l2:crt},
инъективность "--- его единственность (или: множества конечны и равномощны, $\abs{\Zn M} = M =
\prod m_i$). Изоморфизм колец переводит обратимые элементы в обратимые, а элемент прямого
произведения обратим тогда и только тогда, когда обратима каждая его координата; отсюда
изоморфизм групп обратимых элементов и равенство их порядков.
\end{proof}

Изоморфизм~$\varphi$ "--- основа \emph{модулярной арифметики} (системы остаточных классов):
большое число представляется вектором остатков по нескольким небольшим модулям, сложение и
умножение выполняются покомпонентно и независимо (параллельно), а КТО восстанавливает результат.
Тот же приём ускоряет расшифрование RSA (лекция~5): вычисление по модулю $N = pq$ заменяют
двумя вычислениями по модулям $p$ и $q$ вдвое меньшей длины.

\checkyourself{Почему в теореме~\ref{l2:crt} нельзя отказаться от взаимной простоты модулей?
Приведите пример несовместной системы из двух сравнений и пример системы, решение которой не
единственно по модулю $m_1 m_2$.}

\subsection{Алгоритмы и сложность}

\begin{algorithm}[КТО по формуле]\label{l2:alg-crt}
\begin{algorithmic}[1]
\Require попарно взаимно простые $m_1, \dots, m_k$; остатки $r_1, \dots, r_k$.
\Ensure $x \in [0, M)$, $x \equiv r_i \pmod{m_i}$ при всех $i$.
\State $M \gets m_1 m_2 \cdots m_k$, \quad $x \gets 0$
\For{$i = 1, \dots, k$}
  \State $M_i \gets M / m_i$
  \State $y_i \gets (M_i \bmod m_i)^{-1} \bmod m_i$ \Comment{расширенный алгоритм Евклида}
  \State $x \gets (x + r_i M_i y_i) \bmod M$
\EndFor
\State \Return $x$
\end{algorithmic}
\end{algorithm}

Пусть $n = \log_2 M$. Вычисление $M$ и всех $M_i$ требует $\bigO(k)$ умножений и делений
$n$-битного числа на $(\log m_i)$-битное, т.\,е. $\bigO(\sum_i n \log m_i) = \bigO(n^2)$ битовых
операций школьной арифметикой; обращения стоят $\bigO(\sum_i \log^2 m_i) = \bigO(n^2)$;
итоговая сумма "--- ещё $\bigO(n^2)$. Итак, КТО решается за $\bigO(\log^2 M)$ битовых операций.

Другой способ "--- \emph{последовательное решение} (алгоритм Гарнера): присоединяем сравнения по
одному. Будем искать решение в \emph{смешанной системе счисления}
\[
  x = v_1 + v_2 m_1 + v_3 m_1 m_2 + \cdots + v_k m_1 m_2 \cdots m_{k-1}, \qquad 0 \le v_i < m_i.
\]
По модулю $m_1$ все слагаемые, кроме первого, исчезают: $v_1 = r_1 \bmod m_1$. По модулю $m_2$:
$v_1 + v_2 m_1 \equiv r_2$, откуда $v_2 \equiv (r_2 - v_1) m_1^{-1} \pmod{m_2}$, и вообще
\[
  v_i \equiv \bigl(r_i - (v_1 + v_2 m_1 + \cdots + v_{i-1} m_1 \cdots m_{i-2})\bigr)\cdot
  c_i \pmod{m_i}, \qquad c_i = (m_1 \cdots m_{i-1})^{-1} \bmod m_i .
\]

\begin{proposition}\label{l2:garner}
Числа $v_i$ определены однозначно, $x = \sum_i v_i m_1 \cdots m_{i-1}$ удовлетворяет всем
сравнениям системы и $0 \le x < M$.
\end{proposition}

\begin{proof}
Обратные $c_i$ существуют, так как $m_i$ взаимно просто с $m_1 \cdots m_{i-1}$. Слагаемые с номерами
$j > i$ делятся на $m_i$, поэтому $x \equiv v_1 + \cdots + v_i m_1 \cdots m_{i-1} \pmod{m_i}$,
а $v_i$ выбрано ровно так, чтобы правая часть была сравнима с $r_i$. Оценка
$x \le \sum_i (m_i - 1) m_1 \cdots m_{i-1} = M - 1$ получается телескопированием:
$(m_i - 1) m_1 \cdots m_{i-1} = m_1 \cdots m_i - m_1 \cdots m_{i-1}$.
\end{proof}

Достоинства метода Гарнера: при фиксированных модулях константы $c_i$ вычисляются один раз,
а сами вычисления ведутся по \emph{малым} модулям $m_i$ (кроме финальной сборки); кроме того,
смешанное представление $(v_k, \dots, v_1)$ позволяет сравнивать числа по величине, не
восстанавливая их.

\begin{example}
Решим систему примера~\ref{l2:suntzu} методом Гарнера: $v_1 = 2$; $2 + 3v_2 \equiv 3 \pmod 5$,
$v_2 \equiv 1 \cdot 3^{-1} \equiv 2 \pmod 5$, т.\,е. $x \equiv 8 \pmod{15}$; далее
$8 + 15 v_3 \equiv 2 \pmod 7$, $15 \equiv 1$, $v_3 \equiv -6 \equiv 1 \pmod 7$, и
$x = 2 + 2 \cdot 3 + 1 \cdot 15 = 23$.
\end{example}

\checkyourself{Сколько обращений по модулю требуется методу Гарнера и методу «по формуле» для
системы из $k$ сравнений? Какой из них удобнее, если модули фиксированы, а остатки меняются?}

\subsection{Модули, не являющиеся взаимно простыми}

Сначала разберём одно линейное сравнение.

\begin{lemma}\label{l2:linear}
Сравнение $at \equiv b \pmod m$ разрешимо тогда и только тогда, когда $d = \gcd(a, m)$ делит
$b$. В этом случае оно равносильно сравнению $\frac ad\, t \equiv \frac bd \pmod{\frac md}$ и имеет
единственное решение $t_0 \equiv \frac bd \bigl(\frac ad\bigr)^{-1} \pmod{\frac md}$
(т.\,е. ровно $d$ решений по модулю $m$).
\end{lemma}

\begin{proof}
Если $at - b = ms$, то $b = at - ms$ делится на $d$. Обратно, пусть $d \mid b$. Равенство
$at - b = ms$ после деления на $d$ превращается в $\frac ad t - \frac bd = \frac md s$, поэтому
$m \mid at - b \iff \frac md \mid \frac ad t - \frac bd$. Числа $\frac ad$ и $\frac md$ взаимно
просты, и последнее сравнение имеет единственное решение по модулю $\frac md$.
\end{proof}

\begin{theorem}\label{l2:crt-general}
Система $x \equiv r_i \pmod{m_i}$, $i = 1, \dots, k$, совместна тогда и только тогда, когда
\[
  r_i \equiv r_j \pmod{\gcd(m_i, m_j)} \qquad \text{для всех } i < j.
\]
В этом случае решение единственно по модулю $L = \lcm(m_1, \dots, m_k)$.
\end{theorem}

\begin{proof}
\emph{Необходимость.} Если $x$ "--- решение, то $x \equiv r_i$ и $x \equiv r_j$ по модулю
$\gcd(m_i, m_j)$, который делит и $m_i$, и $m_j$; значит, $r_i \equiv r_j \pmod{\gcd(m_i, m_j)}$.

\emph{Достаточность.} По теореме~\ref{l2:crt} (применённой к разложению $m_i = \prod_p p^{a_{i,p}}$)
сравнение $x \equiv r_i \pmod{m_i}$ равносильно системе сравнений
$x \equiv r_i \pmod{p^{a_{i,p}}}$ по всем простым $p \mid m_i$. Зафиксируем простое $p$ и выберем
индекс $i_0$ с наибольшим показателем $a = a_{i_0, p}$. Для любого $j$ число
$p^{\min(a, a_{j,p})} = p^{a_{j,p}}$ делит $\gcd(m_{i_0}, m_j)$, поэтому по условию
$r_{i_0} \equiv r_j \pmod{p^{a_{j,p}}}$: сравнение $x \equiv r_{i_0} \pmod{p^{a}}$ влечёт все
сравнения $x \equiv r_j \pmod{p^{a_{j,p}}}$. Итак, исходная система равносильна системе
\[
  x \equiv r_{i_0(p)} \pmod{p^{a_p}}, \qquad p \mid L, \quad a_p = \max_i a_{i,p},
\]
модули которой попарно взаимно просты, а произведение равно $L$. По КТО она имеет решение,
единственное по модулю $L$.

\emph{Единственность} можно увидеть и напрямую: $x \equiv x'$ по всем модулям $m_i$ тогда и только
тогда, когда $x - x'$ "--- общее кратное всех $m_i$, т.\,е. делится на $L$.
\end{proof}

На практике разложение модулей не нужно: сравнения присоединяют по одному, как в методе Гарнера.
Пусть уже найдено $x \equiv a \pmod{L'}$, и присоединяется $x \equiv r \pmod m$. Подставляя
$x = a + L' t$, получаем линейное сравнение $L' t \equiv r - a \pmod m$; по
лемме~\ref{l2:linear} оно разрешимо тогда и только тогда, когда $d = \gcd(L', m)$ делит $r - a$, и
тогда $t \equiv t_0 \pmod{m/d}$. Новое решение:
\[
  x \equiv a + L' t_0 \pmod{L' m / d}, \qquad L' m / d = \lcm(L', m).
\]
Если $d \nmid r - a$, система несовместна. (Заметим, что $m/d = 1$ означает $m \mid L'$: новое
сравнение либо следует из предыдущих, либо им противоречит.)

\begin{example}\label{l2:noncoprime}
\begin{enumerate}
  \item $x \equiv 3 \pmod 8$, $x \equiv 7 \pmod{12}$. Здесь $\gcd(8, 12) = 4$ и $3 \equiv 7 \pmod 4$
    "--- система совместна. $x = 3 + 8t$, $8t \equiv 4 \pmod{12}$; делим на $4$: $2t \equiv 1 \pmod 3$,
    $t \equiv 2 \pmod 3$. Ответ: $x = 3 + 16 = 19$, $x \equiv 19 \pmod{24}$, $24 = \lcm(8, 12)$.
  \item $x \equiv 1 \pmod 6$, $x \equiv 4 \pmod{10}$: $\gcd(6, 10) = 2$, а $1 - 4 = -3$ нечётно.
    Действительно, первое сравнение требует нечётного $x$, второе "--- чётного: решений нет.
\end{enumerate}
\end{example}

\taskexample{crt-general}

\checkyourself{Сколько решений по модулю $m_1 m_2$ имеет совместная система $x \equiv r_1 \pmod{m_1}$,
$x \equiv r_2 \pmod{m_2}$? Выразите ответ через $\gcd(m_1, m_2)$.}

\section{Логарифмирование в группе порядка $q^e$}

Всюду далее $G$ "--- конечная группа (запись мультипликативная), $g \in G$ имеет порядок $N$,
и решается уравнение $g^x = h$ при $h \in \langle g \rangle$. Ответ определён по модулю $N$.
Нам понадобятся два простых факта о порядках.

\begin{lemma}\label{l2:order}
Пусть $\ord g = N$. Тогда
\begin{enumerate}
  \item $g^a = g^b \iff a \equiv b \pmod N$;
  \item $\ord(g^k) = N / \gcd(N, k)$; в частности, $\ord(g^{N/n}) = n$ для любого делителя $n \mid N$.
\end{enumerate}
\end{lemma}

\begin{proof}
1. $g^a = g^b \iff g^{a-b} = 1 \iff N \mid a - b$ (разделим $a - b$ на $N$ с остатком: остаток
$r < N$ удовлетворяет $g^r = 1$, что по минимальности $N$ возможно лишь при $r = 0$).

2. По п.~1 $(g^k)^s = 1 \iff N \mid ks \iff \frac{N}{d} \mid \frac kd s$, где $d = \gcd(N, k)$;
так как $\frac Nd$ и $\frac kd$ взаимно просты, это равносильно $\frac Nd \mid s$. Наименьшее
такое $s > 0$ равно $N/d$.
\end{proof}

Пусть теперь $N = q^e$, где $q$ "--- простое. Запишем неизвестный показатель в $q$-ичной системе:
\[
  x = x_0 + x_1 q + x_2 q^2 + \cdots + x_{e-1} q^{e-1}, \qquad 0 \le x_i < q .
\]
Элемент $\gamma = g^{q^{e-1}}$ имеет порядок $q$ (лемма~\ref{l2:order}). Ключевое наблюдение:
возведение в степень $q^{e-1}$ «убивает» все цифры, кроме младшей:
\[
  h^{q^{e-1}} = g^{x q^{e-1}} = g^{x_0 q^{e-1}} \cdot \bigl(g^{q^e}\bigr)^{x_1 + x_2 q + \cdots}
  = \gamma^{x_0},
\]
т.\,е. $x_0$ "--- логарифм элемента $h^{q^{e-1}}$ в группе $\langle \gamma \rangle$ порядка $q$.
Найдя $x_0$, «сотрём» младшую цифру и повторим рассуждение для следующей.

\begin{theorem}\label{l2:digits}
Пусть $\ord g = q^e$, $h = g^x$, $x = \sum_{j=0}^{e-1} x_j q^j$, $0 \le x_j < q$,
$\gamma = g^{q^{e-1}}$. Положим
\[
  u_i = h \cdot g^{-(x_0 + x_1 q + \cdots + x_{i-1} q^{i-1})}, \qquad i = 0, 1, \dots, e - 1
\]
($u_0 = h$). Тогда $u_i^{\,q^{e-1-i}} = \gamma^{x_i}$, и цифра $x_i$ однозначно определяется как
логарифм этого элемента по основанию $\gamma$.
\end{theorem}

\begin{proof}
По определению $u_i = g^{x_i q^i + x_{i+1} q^{i+1} + \cdots + x_{e-1} q^{e-1}}$. Возводим в степень
$q^{e-1-i}$:
\[
  u_i^{\,q^{e-1-i}} = g^{x_i q^{e-1}} \cdot g^{q^e (x_{i+1} + x_{i+2} q + \dots)} = \gamma^{x_i},
\]
так как $g^{q^e} = 1$. Элементы $\gamma^0, \gamma^1, \dots, \gamma^{q-1}$ попарно различны
(лемма~\ref{l2:order}, $\ord \gamma = q$), поэтому $x_i \in [0, q)$ находится однозначно.
\end{proof}

Удобно обновлять $u_i$, а не вычислять заново: $u_{i+1} = u_i \cdot (g^{-q^i})^{x_i}$, причём
элементы $c_i = g^{-q^i}$ получаются последовательным возведением $c_0 = g^{-1}$ в степень $q$.
После последнего шага $u_e = h g^{-x} = 1$ "--- удобная проверка.

\begin{algorithm}[DLP в группе порядка $q^e$]\label{l2:alg-digits}
\begin{algorithmic}[1]
\Require $g$ порядка $q^e$, $h \in \langle g \rangle$.
\Ensure $x \in [0, q^e)$, $g^x = h$.
\State $\gamma \gets g^{q^{e-1}}$ \Comment{элемент порядка $q$}
\State $u \gets h$, \quad $c \gets g^{-1}$, \quad $x \gets 0$
\For{$i = 0, 1, \dots, e - 1$}
  \State $w \gets u^{q^{e-1-i}}$ \Comment{$w = \gamma^{x_i}$}
  \State $x_i \gets \log_\gamma w$ \Comment{таблица, Шенкс или $\rho$ в группе порядка $q$}
  \State $x \gets x + x_i q^i$, \quad $u \gets u \cdot c^{x_i}$, \quad $c \gets c^q$
\EndFor
\State \Return $x$
\end{algorithmic}
\end{algorithm}

\begin{proposition}\label{l2:digits-cost}
Пусть DLP в группе порядка $q$ решается за $T_q$ групповых операций. Тогда
алгоритм~\ref{l2:alg-digits} выполняет $\bigO(e \log N + e\,T_q)$ групповых операций, где
$N = q^e$. При $T_q = \bigO(\sqrt q)$ (метод Шенкса) это $\bigO(e(\log N + \sqrt q))$ вместо
$\bigO(\sqrt{N}) = \bigO(q^{e/2})$ для метода Шенкса во всей группе.
\end{proposition}

\begin{proof}
Бинарное возведение в степень $n$ стоит $\bigO(\log n)$ умножений (лекция~6 первого семестра).
На шаге $i$: возведение $u$ в степень $q^{e-1-i} \le N$ "--- $\bigO(\log N)$; умножение на
$c^{x_i}$ и возведение $c$ в степень $q$ "--- $\bigO(\log q) = \bigO(\log N)$; один логарифм
в группе порядка $q$ "--- $T_q$. Вычисление $\gamma$ и $g^{-1} = g^{N-1}$ стоит ещё $\bigO(\log N)$.
Шагов $e$, откуда оценка.
\end{proof}

Если $q$ мало (как в задачах семинара), $\log_\gamma w$ находится по заранее составленной
таблице $\gamma^0, \dots, \gamma^{q-1}$. Особенно прост случай $q = 2$: элемент порядка $2$
в $\Znx p$ "--- это $\gamma = -1$, и бит $x_i$ равен $0$ или $1$ в зависимости от того, равно ли
$u^{2^{e-1-i}}$ единице или $-1$.

\begin{example}\label{l2:ex-digits}
По модулю $p = 433$ элемент $g = 17$ имеет порядок $27 = 3^3$; решим $17^x \equiv 374$.
Здесь $\gamma = 17^9 \equiv 198$, таблица: $\gamma^0 = 1$, $\gamma^1 = 198$, $\gamma^2 \equiv 234$;
$c_0 = 17^{-1} \equiv 51$, $c_1 = 51^3 \equiv 153$.
\[
\begin{array}{c|r|c|r|c|r|r}
  i & u_i & 3^{2-i} & w_i = u_i^{3^{2-i}} & x_i & c_i & u_{i+1} = u_i c_i^{x_i} \\ \hline
  0 & 374 & 9 & 234 & 2 & 51 & 374 \cdot 51^2 \equiv 256 \\
  1 & 256 & 3 & 198 & 1 & 153 & 256 \cdot 153 \equiv 198 \\
  2 & 198 & 1 & 198 & 1 & 234 & 198 \cdot 234 \equiv 1
\end{array}
\]
Итак, $x = 2 + 1 \cdot 3 + 1 \cdot 9 = 14$; действительно, $17^{14} \equiv 374 \pmod{433}$.
\end{example}

\taskexample{ph-prime-power}

\checkyourself{Что произойдёт в алгоритме~\ref{l2:alg-digits}, если $h = 1$? Если $h = g^{-1}$?
Если $h$ имеет порядок $q$?}

\checkyourself{Почему цифры $x$ находятся начиная с \emph{младшей}? Можно ли найти сначала старшую
цифру $x_{e-1}$?}

\section{Алгоритм Полига"--~Хеллмана}

Теперь пусть $\ord g = N = q_1^{e_1} q_2^{e_2} \cdots q_t^{e_t}$ "--- произвольное число с известным
разложением на простые множители.

\begin{theorem}[Полиг"--~Хеллман~\cite{pohlig-hellman}]\label{l2:ph}
Пусть $\ord g = N = \prod_{i=1}^{t} q_i^{e_i}$ и $h \in \langle g \rangle$. Для каждого $i$ положим
\[
  g_i = g^{N / q_i^{e_i}}, \qquad h_i = h^{N / q_i^{e_i}} ,
\]
и пусть $y_i$ "--- решение уравнения $g_i^{\,y} = h_i$. Тогда $\ord g_i = q_i^{e_i}$, и решение $x$
системы
\[
  x \equiv y_1 \pmod{q_1^{e_1}}, \quad \dots, \quad x \equiv y_t \pmod{q_t^{e_t}}
\]
(оно существует и единственно по модулю $N$ по КТО) удовлетворяет равенству $g^x = h$.
\end{theorem}

\begin{proof}
По лемме~\ref{l2:order} $\ord g_i = q_i^{e_i}$. Так как $h \in \langle g \rangle$, имеем $h = g^{z}$
для некоторого $z$, и тогда $h_i = g^{z N / q_i^{e_i}} = g_i^{\,z}$. Значит, $g_i^{\,y_i} = g_i^{\,z}$,
и по лемме~\ref{l2:order} $y_i \equiv z \pmod{q_i^{e_i}}$. Числа $x$ и $z$ удовлетворяют одной и той же
системе сравнений по попарно взаимно простым модулям $q_i^{e_i}$, поэтому по теореме~\ref{l2:crt}
$x \equiv z \pmod N$ и $g^x = g^z = h$.
\end{proof}

\begin{remark}
Условие $h \in \langle g \rangle$ существенно: иначе уравнение не имеет решения, а алгоритм выдаст
какое-то $x$. В циклической группе (например, $\Znx p$) подгруппа порядка $N$ единственна и
состоит в точности из элементов с $h^N = 1$, так что условие проверяется одним возведением в
степень. В любом случае ответ стоит проверить: $g^x = h$.
\end{remark}

\begin{algorithm}[Полиг"--~Хеллман]\label{l2:alg-ph}
\begin{algorithmic}[1]
\Require $g$ порядка $N = \prod_{i=1}^{t} q_i^{e_i}$ (разложение известно), $h \in \langle g \rangle$.
\Ensure $x \in [0, N)$, $g^x = h$.
\For{$i = 1, \dots, t$}
  \State $g_i \gets g^{N / q_i^{e_i}}$, \quad $h_i \gets h^{N / q_i^{e_i}}$
  \State $y_i \gets \log_{g_i} h_i$ \Comment{алгоритм~\ref{l2:alg-digits}, группа порядка $q_i^{e_i}$}
\EndFor
\State найти $x$: $x \equiv y_i \pmod{q_i^{e_i}}$, $i = 1, \dots, t$ \Comment{КТО}
\State \Return $x$
\end{algorithmic}
\end{algorithm}

\begin{theorem}\label{l2:ph-cost}
Алгоритм Полига"--~Хеллмана, использующий для подгрупп простого порядка метод Шенкса, решает DLP
в группе порядка $N = \prod_{i=1}^{t} q_i^{e_i}$ за
\[
  \bigO\Bigl(\sum_{i=1}^{t} e_i \bigl(\log N + \sqrt{q_i}\bigr)\Bigr)
\]
групповых операций и $\bigO(\log^2 N)$ битовых операций на сборку ответа по КТО.
\end{theorem}

\begin{proof}
Для каждого $i$ вычисление $g_i$ и $h_i$ "--- два возведения в степень, не превосходящую $N$:
$\bigO(\log N)$ операций, что не больше $\bigO(e_i \log N)$. Задача в подгруппе порядка
$q_i^{e_i}$ решается алгоритмом~\ref{l2:alg-digits} за $\bigO(e_i \log N + e_i \sqrt{q_i})$
операций (утверждение~\ref{l2:digits-cost}). Суммируя по $i$, получаем первое слагаемое.
Решение системы по КТО с модулем $N$ стоит $\bigO(\log^2 N)$ битовых операций (раздел о КТО).
\end{proof}

Поскольку $\sum_i e_i \le \log_2 N$, оценку можно огрубить до
$\bigO\bigl(\log^2 N + \log N \cdot \sqrt{q_{\max}}\bigr)$, где $q_{\max}$ "--- наибольший простой
делитель $N$. Иначе говоря, \emph{сложность DLP в группе порядка $N$ не превосходит сложности DLP
в её подгруппе наибольшего простого порядка} (с полиномиальным множителем). Если $N$ "---
$B$-гладкое (все простые делители не больше $B$) и $B$ мало, логарифм вычисляется практически
мгновенно.

\begin{example}\label{l2:ex-ph}
$p = 433$, $N = p - 1 = 432 = 2^4 \cdot 3^3$, $g = 5$ "--- первообразный корень; решим
$5^x \equiv 7 \pmod{433}$.

\emph{Подгруппа порядка $16$:} $g_1 = 5^{27} \equiv 238$, $h_1 = 7^{27} \equiv 265$,
$\gamma = 238^8 \equiv 432 = -1$, $c_0 = 238^{-1} \equiv 282$.
\[
\begin{array}{c|r|c|r|c|r|r}
  i & u_i & 2^{3-i} & w_i & x_i & c_i & u_{i+1} \\ \hline
  0 & 265 & 8 & -1 & 1 & 282 & 254 \\
  1 & 254 & 4 & 1 & 0 & 285 & 254 \\
  2 & 254 & 2 & -1 & 1 & 254 & 432 \\
  3 & 432 & 1 & -1 & 1 & 432 & 1
\end{array}
\qquad y_1 = 1 + 4 + 8 = 13 .
\]

\emph{Подгруппа порядка $27$:} $g_2 = 5^{16} \equiv 17$, $h_2 = 7^{16} \equiv 374$ "--- это в точности
пример~\ref{l2:ex-digits}: $y_2 = 14$.

\emph{Сборка по КТО:} $x \equiv 13 \pmod{16}$, $x \equiv 14 \pmod{27}$. Методом Гарнера:
$x = 13 + 16t$, $16t \equiv 1 \pmod{27}$, $16^{-1} \equiv 22 \pmod{27}$, $t \equiv 22$, и
$x = 13 + 16 \cdot 22 = 365$. Проверка: $5^{365} \equiv 7 \pmod{433}$.

Вместо $\sqrt{432} \approx 21$ шагов метода Шенкса мы решили $4 + 3$ «задачи» в группах порядков
$2$ и $3$, т.\,е. просто сравнивали элементы с таблицами из двух и трёх элементов.
\end{example}

В следующем примере (сгенерированном taskgen) частные задачи в подгруппах решаются перебором
степеней $g_i$ "--- при малых $q_i^{e_i}$ это тоже разумный способ.

\taskexample{ph-full}

\checkyourself{Сколько групповых операций потребует алгоритм Полига"--~Хеллмана в $\Znx p$ для
простого Ферма $p = 2^{16} + 1$? А для $p = 2q + 1$, где $q$ "--- простое число из $255$ бит?}

\section{Следствия для криптографии}

Теорема~\ref{l2:ph-cost} диктует выбор параметров всех протоколов, основанных на DLP.

\begin{itemize}
  \item \textbf{Порядок группы должен иметь большой простой делитель.} Для уровня стойкости
    $\lambda$ бит нужна подгруппа простого порядка $q \ge 2^{2\lambda}$, так как универсальные
    алгоритмы (Шенкс, $\rho$-метод Полларда "--- лекция~3) тратят $\sqrt q$ операций. Для
    $\lambda = 128$ это $q$ длиной около $256$ бит.
  \item \textbf{Гладкое $p - 1$ "--- слабость.} Если все простые делители $p - 1$ малы (например,
    $p = 2^{16} + 1$ или $p = 2^a 3^b + 1$), DLP в $\Znx p$ решается за полиномиальное время,
    какой бы большой ни была длина $p$. Протокол Диффи"--~Хеллмана с таким $p$ не стоек.
  \item \textbf{Безопасные простые} $p = 2q + 1$, $q$ "--- простое. Тогда $p - 1 = 2q$, и
    алгоритм Полига"--~Хеллмана ничего не выигрывает. Однако он всё же извлекает \emph{один бит}:
    $x \bmod 2$ определяется тем, является ли $h$ квадратичным вычетом ($h^{(p-1)/2} = \pm 1$).
    Поэтому обычно работают в подгруппе квадратичных вычетов порядка $q$.
  \item \textbf{Группы Шнорра}: $p = kq + 1$, где $q$ "--- простое длиной $256$ бит, $p$ "--- длиной
    $3072$ бит (такой размер $p$ нужен из-за субэкспоненциального метода исчисления индексов,
    лекция~4). Рабочая подгруппа порядка $q$ порождается элементом $g = a^{(p-1)/q} \ne 1$;
    показатели берутся по модулю $q$ и потому короткие.
  \item \textbf{Атаки малых подгрупп.} Если участник протокола возводит в свою секретную степень $x$
    элемент, присланный противником, и не проверяет, что тот лежит в подгруппе порядка $q$,
    противник может прислать элемент малого порядка $r \mid p - 1$ и по ответу узнать $x \bmod r$;
    собрав такие остатки по КТО, он восстановит $x$. Защита "--- проверка $h^q = 1$.
  \item Те же соображения относятся к любой группе, в частности к группе точек эллиптической
    кривой (лекция~9): её порядок должен делиться на большое простое число.
\end{itemize}

\checkyourself{Почему для работы алгоритма Полига"--~Хеллмана нужно знать разложение $N$ на
множители? Почему в $\Znx p$ это обычно не проблема, хотя факторизация "--- трудная задача?}

\checkyourselfsection

\begin{furtherreading}
  \readitem[разд.~2.8, 2.9]{hps}{китайская теорема об остатках, алгоритм Полига"--~Хеллмана
    (основной источник лекции), упражнения к этим разделам}
  \readitem[разд.~13.2]{galbraith}{алгоритм Полига"--~Хеллмана в произвольной группе, оценки
    сложности}
  \readitem[разд.~3.6.4, 14.5]{hac}{алгоритм Полига"--~Хеллмана; КТО и алгоритм Гарнера}
  \readitem[разд.~4.3.2]{knuth-taocp2}{модулярная арифметика, КТО, смешанная система счисления
    и метод Гарнера}
  \readitem[разд.~2.1.3]{crandall-pomerance}{китайская теорема об остатках и её алгоритмы}
  \readitem[разд.~5.4]{von-zur-gathen}{китайский алгоритм остатков и его сложность}
  \readitem{pohlig-hellman}{оригинальная статья}
  \readitem{cp-algorithms}{статьи «Chinese Remainder Theorem» и «Garner's algorithm» с
    реализациями}
  \readitem{mit-18783}{конспект «Generic algorithms for the discrete logarithm problem»:
    алгоритм Полига"--~Хеллмана для абстрактной группы}
\end{furtherreading}
\end{lecture}
